
\begin{frame}[fragile]{Le modèle relationnel}

Le modèle relationnel, c'est

\begin{redbullet}
  \item Une structure unique, la relation (ou table)
  \item Des contraintes qui définissent des \alert{formes normales}, évitant les défauts de conception
  \item Des langages, concrétisés en pratique par SQL
\end{redbullet}

\vfill

Dans ce cours, on parle de la structure

\vfill
\begin{blocgauche}
\alert{Ces diapositives correspondent au support en ligne disponible sur
le site \url{http://sql.bdpedia.fr/}}
\end{blocgauche}

\end{frame}


\begin{frame}{Qu'est-ce qu'une relation?}

\alert{Notion mathématique}: 
Etant donné un ensemble d'objets $O$, une \alert{relation} (binaire)
sur $O$  est un sous-ensemble du produit cartésien $O \times O$.

\vfill
 Dans notre contexte, les ``objets''  sont des \alert{valeurs élémentaires} (ou \alert{atomiques}), comme 
 les entiers $I$, les flottants $F$, les chaînes de caractères $S$.


\vfill
\begin{blocgauche}
L'ensemble des paires constituées des noms de département et 
et de leur numéro de code  est une relation sur $S \times I$. 

\begin{defblock}{Définition: relation}
Une relation de degré $n$ sur les domaines $A_1, A_2, \cdots, A_n$
    est un sous-ensemble fini du produit cartésien  $A_1 \times A_2 \times  \cdots \times A_n$
\end{defblock}
\end{blocgauche}
\end{frame}


\begin{frame}{Représentation}

Comment représente-t-on une relation? Sous forme de table, le plus pratique.

\vfill

\begin{tabular}{l|l}
\textbf{nom} & \textbf{code}\\ \hline
  Ardèche & 07\\
   Gard & 30\\
   Manche & 50\\
   Paris& 75\\ \hline
\end{tabular}

\vfill

\begin{blocgauche}
Attention, ce n'est pas \alert{n'importe quelle table}. Se
souvenir de la définition.
\end{blocgauche}

\end{frame}


%%%%%%%%%%%%%%%
\begin{frame}{Les nuplets}

Un élément d'une relation de dimension $n$ est un \alert{nuplet}
$(a_1, a_2, \cdots, a_n)$.

\vfill

Exemple: \texttt{(Manche, 50)}
\vfill

Dans la représentation par table, 
un nuplet est une ligne. 

\vfill

\begin{blocgauche}
On assimile nuplet et ligne, mais attention, ce n'est pas \alert{n'importe quelle ligne}. Se
souvenir de la définition.
\end{blocgauche}

\end{frame}


%%%%%%%%%%%%%%%
\begin{frame}{Schéma de relation}

On peut  \alert{décrire} une relation par 

\begin{redbullet}
  \item Le nom de la relation.
  \item Un nom (distinct) pour chaque dimension, dit \alert{nom d'attribut}
  \item Le domaine de valeur de chaque dimension.
\end{redbullet}

\vfill

C'est le \alert{schéma} de la relation, de la forme 
$R (A_1: D_1, A_2: D_2, \cdots, A_n: D_n)$

\vfill
\begin{blocgauche}
\alert{Exemple}:  \texttt{Département (nom: string, code: string)}, ou plus simplement
\texttt{Département (nom, code)}
\end{blocgauche}
\end{frame}


%%%%%%%%%%%%%%%
\begin{frame}{Première forme normale}

 On \alert{ne peut
pas} avoir une valeur d'attribut qui soit construite, comme par
exemple une liste, ou une sous-relation. 

\vfill

Les valeurs dans une base de données sont dites \alert{atomiques} 

\vfill

\begin{defblock}{Première forme normale}
Une relation est en première forme normale si toutes les valeurs d'attribut sont connues et atomiques
et si elle ne contient aucun doublon.
\end{defblock}

\end{frame}



\begin{frame}{À retenir}

\alert{Notions et vocabulaire}

\vfill

\begin{tabular}{l|l}
\textbf{Terme du modèle} & \textbf{Terme de la représentation par table}\\ \hline
 Relation& Table\\
   nuplet& ligne\\
   Nom d'attribut &Nom de colonne\\
   Valeur d'attribut& Cellule\\
   Domaine&Type\\
\hline
\end{tabular}

\vfill

\alert{Très simple\,!} (parfois trop). Favorise la rigueur et la clarté.
\end{frame}

